\setcounter{section}{6}

  \zwischenueberschrift{Soal latihan}

\inputaufgabe
{}
{

Misalkan $Y$ adalah
\definitionsverweis {grafik}{}{}
fungsi

\mavergleichskette
{\vergleichskette
{ f(x,y)
}
{ = }{ x^2-y^2
}
{  }{
}
{    }{
}
{   }{
}
}
{}{}{.}
\aufzaehlungdrei{Buktikan bahwa

\mavergleichskettedisp
{\vergleichskette
{ F(x,y,z)
}
{ =} { \begin{pmatrix}  y  \\ x \\  0   \end{pmatrix}
}
{ } {
}
{ } {
}
{ } {
}
}
{}{}{}
memberikan suatu
\definitionsverweis {medan vektor tangensial}{}{}
pada $Y$.
}{Misalkan

\mavergleichskette
{\vergleichskette
{ P
}
{ = }{ \begin{pmatrix}  2  \\ 3 \\  -5   \end{pmatrix}
}
{ \in }{ Y
}
{    }{
}
{   }{
}
}
{}{}{.}
Hitung
\mathl{\nabla_v F}{} untuk

\mavergleichskette
{\vergleichskette
{ v
}
{ = }{ \begin{pmatrix}  2  \\ 1 \\  2   \end{pmatrix}
}
{ \in }{ T_PY
}
{    }{
}
{   }{
}
}
{}{}{.}
\emph{Catatan edisi: Tanda sama dengan pada sumber diperbaiki menjadi tanda keanggotaan, karena $v$ adalah vektor di $T_PY$.}
}{Verifikasikan bahwa $\nabla_v F$ tangensial pada $Y$ di $P$.
}

}
{} {}

\inputaufgabegibtloesung
{}
{

Pada
\definitionsverweis {permukaan bola satuan}{}{}
$Y$ yang berdimensi dua, kita tinjau
\definitionsverweis {medan vektor tangensial}{}{}

\mavergleichskettedisp
{\vergleichskette
{ F \begin{pmatrix}  x  \\ y \\ z   \end{pmatrix}
}
{ =} { \begin{pmatrix}  y  \\ -x\\  0   \end{pmatrix}
}
{ } {
}
{ } {
}
{ } {
}
}
{}{}{.}
\aufzaehlungdrei{Tentukan
\definitionsverweis {matriks Jacobi}{}{}
dari $F$ pada $\R^3$.
}{Untuk titik

\mavergleichskette
{\vergleichskette
{ P
}
{ = }{ \begin{pmatrix}  0  \\ 1 \\  0   \end{pmatrix}
}
{ \in }{ Y
}
{    }{
}
{   }{
}
}
{}{}{}
tentukan suatu matriks bagi pemetaan
\maabbeledisp {} {T_PY} { T_PY
} { v } { \nabla_v F
} {,}
terhadap suatu
\definitionsverweis {basis}{}{}
yang sesuai.
}{Untuk titik

\mavergleichskette
{\vergleichskette
{ Q
}
{ = }{ \begin{pmatrix}  0  \\ 0\\  1   \end{pmatrix}
}
{ \in }{ Y
}
{    }{
}
{   }{
}
}
{}{}{}
tentukan suatu matriks bagi pemetaan
\maabbeledisp {} {T_QY} { T_QY
} { v } { \nabla_v F
} {,}
terhadap suatu
\definitionsverweis {basis}{}{}
yang sesuai.
}

}
{} {}

\inputaufgabe
{}
{

Misalkan

\mavergleichskette
{\vergleichskette
{ Y
}
{ \subseteq }{ W
}
{ \subseteq }{ \R^n
}
{    }{
}
{   }{
}
}
{}{}{,}
dengan $W$
\definitionsverweis {terbuka}{}{,}
adalah suatu
\definitionsverweis {hipermuka terdiferensialkan}{}{,}
dan misalkan
\maabb {\gamma} { I } { Y
} {}
adalah suatu
\definitionsverweis {kurva terdiferensialkan}{}{.}
Buktikan pernyataan-pernyataan berikut.
\aufzaehlungzwei {Jika $F$ adalah suatu
\definitionsverweis {medan vektor paralel}{}{}
sepanjang $\gamma$ dan

\mavergleichskette
{\vergleichskette
{ a
}
{ \in }{ \R
}
{  }{
}
{    }{
}
{   }{
}
}
{}{}{,}
maka \mathl{aF}{} juga merupakan medan vektor paralel.
} {Jika $F$ dan $G$ adalah medan-medan vektor paralel sepanjang $\gamma$, maka
\mathl{F +G}{} juga merupakan medan vektor paralel.
}

}
{} {}

\inputaufgabe
{}
{

Misalkan

\mavergleichskette
{\vergleichskette
{ Y
}
{ \subseteq }{ \R^n
}
{  }{
}
{    }{
}
{   }{
}
}
{}{}{}
adalah suatu
\definitionsverweis {hiperbidang}{}{.}
Buktikan bahwa setiap
\definitionsverweis {medan vektor tangensial}{}{}
$F$ yang konstan
\zusatzklammer {yang bersesuaian dengan suatu medan vektor pada

\mavergleichskette
{\vergleichskette
{ Y
}
{ \cong }{\R^{n-1}
}
{  }{
}
{    }{
}
{   }{
}
}
{}{}{}} {} {}
bersifat
\definitionsverweis {paralel}{}{}
\zusatzklammer {sepanjang setiap kurva} {}{}{.}

\emph{Catatan edisi: Syarat ``konstan'' ditambahkan; medan tangensial umum pada hiperbidang tidak harus paralel.}

}
{} {}

\inputaufgabe
{}
{

Misalkan

\mavergleichskettedisp
{\vergleichskette
{ Y
}
{ =} { \left\{  \begin{pmatrix}  x  \\ y \\ z   \end{pmatrix}   \mid   x^2+y^2+z^2 = 1        \right\}
}
{ } {
}
{ } {
}
{ } {
}
}
{}{}{}
adalah
\definitionsverweis {permukaan bola satuan}{}{}
berdimensi dua dan
\maabbele {\gamma} {[0,2\pi]} {Y
} { t } { \begin{pmatrix}  \cos  t  \\ 0 \\  \sin  t   \end{pmatrix}
} {.}
Tentukan mana di antara medan-medan vektor berikut
\maabbdisp {F} {[0,2\pi]} { \R^3
} {}
sepanjang $\gamma$ yang tangensial pada $Y$ dan mana yang
\definitionsverweis {paralel}{}{.}
\aufzaehlungfuenf{
\mavergleichskette
{\vergleichskette
{ F (t)
}
{ = }{ \begin{pmatrix}  2  \\ 0 \\  1   \end{pmatrix}
}
{  }{
}
{    }{
}
{   }{
}
}
{}{}{.}
}{
\mavergleichskette
{\vergleichskette
{ F (t)
}
{ = }{ \begin{pmatrix}  0  \\ 5 \\  0   \end{pmatrix}
}
{  }{
}
{    }{
}
{   }{
}
}
{}{}{.}
}{
\mavergleichskette
{\vergleichskette
{ F (t)
}
{ = }{ \begin{pmatrix}  \cos  t  \\ 0 \\  \sin  t   \end{pmatrix}
}
{  }{
}
{    }{
}
{   }{
}
}
{}{}{.}
}{
\mavergleichskette
{\vergleichskette
{ F (t)
}
{ = }{ \begin{pmatrix}  - \sin  t  \\ 0 \\  \cos  t   \end{pmatrix}
}
{  }{
}
{    }{
}
{   }{
}
}
{}{}{.}
}{
\mavergleichskette
{\vergleichskette
{ F (t)
}
{ = }{ \begin{pmatrix}  t  \\e^t\\  \sin  t   \end{pmatrix}
}
{  }{
}
{    }{
}
{   }{
}
}
{}{}{.}
}

}
{} {}

\inputaufgabegibtloesung
{}
{

Misalkan

\mavergleichskette
{\vergleichskette
{ Y
}
{ \subseteq }{ W
}
{ \subseteq }{ \R^3
}
{    }{
}
{   }{
}
}
{}{}{,}
dengan $W$
\definitionsverweis {terbuka}{}{,}
adalah suatu
\definitionsverweis {permukaan terdiferensialkan}{}{}
yang dilengkapi dengan
\definitionsverweis {medan normal satuan}{}{}
$N$. Misalkan
\maabb {\gamma} { I } { Y
} {}
adalah suatu kurva yang terdiferensialkan secara kontinu dan
\maabbdisp {F} { I } { \R^3
} {}
adalah suatu
\definitionsverweis {medan vektor tangensial}{}{}
terdiferensialkan sepanjang $\gamma$ yang
\definitionsverweis {paralel}{}{.}
Buktikan bahwa medan vektor
\maabbeledisp {} { I } { \R^3
} { t } { N(\gamma(t)) \times F(t)
} {,}
juga paralel, dengan $\times$ menyatakan
\definitionsverweis {hasil kali silang}{}{}
di $\R^3$.

}
{} {}

\inputaufgabe
{}
{

Misalkan

\mavergleichskettedisp
{\vergleichskette
{ h(x,y,z)
}
{ =} { x^3-2y^4-3z^5
}
{ } {
}
{ } {
}
{ } {
}
}
{}{}{}
dan

\mavergleichskette
{\vergleichskette
{ Y
}
{ = }{ h^{-1}(1)
}
{  }{
}
{    }{
}
{   }{
}
}
{}{}{.}
Misalkan
\maabbeledisp {\gamma} { (-3^{-1/5},\infty)} {Y
} { t } { \begin{pmatrix}  \sqrt[3]{ 1+3t^5  } \\ 0 \\ t   \end{pmatrix}
} {,}
adalah suatu lintasan terdiferensialkan pada $Y$. Tentukan persamaan diferensial biasa dalam bentuk
[[Hyperfläche/Kurve/Anfangstangentenvektor/Paralleles Vektorfeld/Differentialgleichung/Explizite Gestalt/Bemerkung|Catatan 6.10]]
untuk
\definitionsverweis {medan-medan vektor paralel}{}{}
sepanjang $\gamma$.

\emph{Catatan edisi: Domain sumber $\R$ dibatasi pada interval di atas karena komponen akar pangkat tiga tidak terdiferensialkan di $t=-3^{-1/5}$.}

}
{} {}

\inputaufgabe
{}
{

Misalkan

\mavergleichskette
{\vergleichskette
{ Y
}
{ \subseteq }{ \R^n
}
{  }{
}
{    }{
}
{   }{
}
}
{}{}{}
adalah suatu
\definitionsverweis {hiperbidang}{}{.}
Buktikan bahwa
\definitionsverweis {transpor paralel}{}{}
sepanjang setiap kurva pada $Y$ adalah identitas setelah semua
\definitionsverweis {ruang singgung}{}{}
$T_PY$ diidentifikasi secara kanonis dengan ruang arah linear $V=Y-P$ dari hiperbidang tersebut.

\emph{Catatan edisi: Sumber mengidentifikasi ruang singgung hiperbidang afin langsung dengan $Y$. Identifikasi kanonis sebenarnya adalah dengan ruang arah linear $V=Y-P$; pernyataan diperjelas sesuai identifikasi tersebut.}

}
{} {}

\inputaufgabegibtloesung
{}
{

Misalkan

\mavergleichskette
{\vergleichskette
{ Y
}
{ \subseteq }{ W
}
{ \subseteq }{ \R^n
}
{    }{
}
{   }{
}
}
{}{}{,}
dengan $W$
\definitionsverweis {terbuka}{}{,}
adalah suatu
\definitionsverweis {hipermuka terdiferensialkan}{}{,}
dan misalkan
\maabb {\gamma} {I =[a,b]} {Y
} {}
adalah suatu
\definitionsverweis {kurva geodesik}{}{.}
Buktikan bahwa
\definitionsverweis {transpor paralel}{}{}
sepanjang $\gamma$ membawa vektor singgung

\mavergleichskette
{\vergleichskette
{ \gamma'(a)
}
{ \in }{ T_{\gamma(a)}Y
}
{  }{
}
{    }{
}
{   }{
}
}
{}{}{}
ke vektor singgung

\mavergleichskette
{\vergleichskette
{ \gamma'(b)
}
{ \in }{ T_{\gamma(b)}Y
}
{  }{
}
{    }{
}
{   }{
}
}
{}{}{.}

}
{} {}

\inputaufgabe
{}
{

Misalkan $Y$ adalah
\definitionsverweis {permukaan bola satuan}{}{}
berdimensi dua dan
\maabb {\gamma} {[0,2\pi]} {Y
} {}
adalah
\definitionsverweis {parametrisasi panjang busur}{}{}
dari suatu
\definitionsverweis {lingkaran besar}{}{}
dengan

\mavergleichskette
{\vergleichskette
{ \gamma(0)
}
{ = }{ \gamma(2 \pi)
}
{ = }{ P
}
{ \in   }{ Y
}
{   }{
}
}
{}{}{.}
Buktikan bahwa
\definitionsverweis {transpor paralel}{}{}
yang bersesuaian adalah identitas pada $T_PY$.

}
{} {}

\inputaufgabe
{}
{

Misalkan

\mavergleichskettedisp
{\vergleichskette
{ Y
}
{ =} { \left\{  \begin{pmatrix}  x  \\ y \\ z   \end{pmatrix}   \mid   x^2+y^2+z^2 = 1        \right\}
}
{ } {
}
{ } {
}
{ } {
}
}
{}{}{}
adalah
\definitionsverweis {permukaan bola satuan}{}{}
berdimensi dua,

\mavergleichskette
{\vergleichskette
{ P
}
{ = }{ \begin{pmatrix}  { \frac{   1  }{  2 } }  \\ 0 \\  { \frac{   \sqrt{3}  }{  2 } }   \end{pmatrix}
}
{  }{
}
{    }{
}
{   }{
}
}
{}{}{}
dan
\maabbele {\gamma} {[0,2\pi]} {Y
} { t } { \begin{pmatrix}  { \frac{   1  }{  2  } } \cos  t  \\ { \frac{   1  }{  2  } }  \sin  t \\  { \frac{   \sqrt{3}  }{  2  } }   \end{pmatrix}
} {.}
Deskripsikan
\definitionsverweis {transpor paralel}{}{}
\maabb {\Psi_\gamma} { T_PY} {T_PY
} {}
terhadap suatu basis yang sesuai.

}
{} {}

\inputaufgabe
{}
{

Misalkan

\mavergleichskette
{\vergleichskette
{ Y
}
{ \subseteq }{ W
}
{ \subseteq }{ \R^n
}
{    }{
}
{   }{
}
}
{}{}{,}
dengan $W$
\definitionsverweis {terbuka}{}{,}
adalah suatu
\definitionsverweis {hipermuka terdiferensialkan}{}{,}
dan misalkan

\mavergleichskette
{\vergleichskette
{ P,Q
}
{ \in }{ Y
}
{  }{
}
{    }{
}
{   }{
}
}
{}{}{}
adalah titik-titik dengan
\definitionsverweis {ruang-ruang singgung}{}{}
\mathkor {} {T_PY} {dan} {T_QY} {.}
Kita tinjau pemetaan
\maabbeledisp {\varphi_{P,Q}} {T_PY} { T_QY
} { v } { \pi_Q(v)
} {,}
yang memandang suatu vektor singgung di $P$ sebagai vektor di ruang sekitar $\R^n$, lalu memproyeksikannya secara ortogonal ke $T_QY$.
\aufzaehlungvier{Apakah
\mathl{\varphi_{P,Q}}{} linear?
}{Apakah
\mathl{\varphi_{P,Q}}{} bijektif?
}{Apakah
\mathl{\varphi_{P,Q}}{} suatu isometri?
}{Apakah terdapat hubungan antara
\mathl{\varphi_{P,Q}}{} dan
\definitionsverweis {transpor paralel}{}{} $\Psi_\gamma$ sepanjang suatu kurva terdiferensialkan
\maabb {\gamma} {[a,b]} { Y
} {}
dengan

\mavergleichskette
{\vergleichskette
{ \gamma(a)
}
{ = }{ P
}
{  }{
}
{    }{
}
{   }{
}
}
{}{}{}
dan

\mavergleichskette
{\vergleichskette
{ \gamma(b)
}
{ = }{ Q
}
{  }{
}
{    }{
}
{   }{
}
}
{}{}{?}
}

}
{} {}

Misalkan

\mavergleichskette
{\vergleichskette
{ W
}
{ \subseteq }{ \R^n
}
{  }{
}
{    }{
}
{   }{
}
}
{}{}{}
\definitionsverweis {terbuka}{}{,}
\maabb {h} { W } { \R
} {}
adalah suatu
\definitionsverweis {fungsi yang dua kali terdiferensialkan secara kontinu}{}{}
dan

\mavergleichskette
{\vergleichskette
{ Y
}
{ = }{ h^{-1}(c)
}
{  }{
}
{    }{
}
{   }{
}
}
{}{}{}
adalah
\definitionsverweis {serat}{}{}
di atas

\mavergleichskette
{\vergleichskette
{ c
}
{ \in }{ \R
}
{  }{
}
{    }{
}
{   }{
}
}
{}{}{,}
dengan $h$
\definitionsverweis {reguler}{}{}
di setiap titik pada $Y$. Untuk suatu titik

\mavergleichskette
{\vergleichskette
{ P
}
{ \in }{ Y
}
{  }{
}
{    }{
}
{   }{
}
}
{}{}{,}
\definitionsverweis {subgrup}{}{}

\mavergleichskette
{\vergleichskette
{ H
}
{ \subseteq }{ \operatorname{Iso} \left(  T_PY  \right)
}
{  }{
}
{    }{
}
{   }{
}
}
{}{}{,}
yang terdiri atas semua
\definitionsverweis {transpor paralel}{}{}
\maabbdisp {\Psi_\gamma} {T_PY} {T_PY
} {}
yang diinduksi oleh
\definitionsverweis {lintasan tertutup}{}{}
$\gamma$ yang terdiferensialkan sepotong-sepotong pada $Y$ dari $P$ ke $P$, disebut
\definitionswort {grup holonomi}{}
di $P$.

\emph{Catatan edisi: Regularitas sumber diperkuat dari $C^1$ menjadi $C^2$ agar medan normal dapat diturunkan dan transpor paralel yang dipakai dalam definisi ini terjamin.}

\inputaufgabe
{}
{

Misalkan

\mavergleichskette
{\vergleichskette
{ W
}
{ \subseteq }{ \R^n
}
{  }{
}
{    }{
}
{   }{
}
}
{}{}{}
\definitionsverweis {terbuka}{}{,}
\maabb {h} { W } { \R
} {}
adalah suatu
\definitionsverweis {fungsi yang dua kali terdiferensialkan secara kontinu}{}{}
dan

\mavergleichskette
{\vergleichskette
{ Y
}
{ = }{ h^{-1}(c)
}
{  }{
}
{    }{
}
{   }{
}
}
{}{}{}
adalah
\definitionsverweis {serat}{}{}
di atas

\mavergleichskette
{\vergleichskette
{ c
}
{ \in }{ \R
}
{  }{
}
{    }{
}
{   }{
}
}
{}{}{,}
dengan $h$
\definitionsverweis {reguler}{}{}
di setiap titik pada $Y$. Buktikan bahwa
\definitionsverweis {grup holonomi}{}{}
di suatu titik

\mavergleichskette
{\vergleichskette
{ P
}
{ \in }{ Y
}
{  }{
}
{    }{
}
{   }{
}
}
{}{}{}
memang merupakan suatu
\definitionsverweis {subgrup}{}{}
dari
\definitionsverweis {grup isometri}{}{.}

\emph{Catatan edisi: Hipotesis $C^1$ pada sumber diperkuat menjadi $C^2$, sesuai dengan regularitas yang diperlukan oleh transpor paralel.}

}
{} {}

\inputaufgabe
{}
{

Misalkan

\mavergleichskette
{\vergleichskette
{ Y
}
{ \subseteq }{ \R^n
}
{  }{
}
{    }{
}
{   }{
}
}
{}{}{}
adalah suatu
\definitionsverweis {hiperbidang}{}{.}
Tentukan
\definitionsverweis {grup holonomi}{}{}
di

\mavergleichskette
{\vergleichskette
{ P
}
{ \in }{ Y
}
{  }{
}
{    }{
}
{   }{
}
}
{}{}{.}

}
{} {}

  \zwischenueberschrift{Soal untuk dikumpulkan}

\inputaufgabe
{2}
{

Pada
\definitionsverweis {permukaan bola satuan}{}{}
$Y$ yang berdimensi dua, kita tinjau
\definitionsverweis {medan vektor tangensial}{}{}

\mavergleichskettedisp
{\vergleichskette
{ F \begin{pmatrix}  x  \\ y \\ z   \end{pmatrix}
}
{ =} { \begin{pmatrix}  y  \\ -x\\  0   \end{pmatrix}
}
{ } {
}
{ } {
}
{ } {
}
}
{}{}{.}
Untuk titik

\mavergleichskette
{\vergleichskette
{ R
}
{ = }{ \begin{pmatrix}  { \frac{   1  }{  2 } }  \\ { \frac{   1  }{  2 } } \\  { \frac{   1  }{  \sqrt{2}  } }   \end{pmatrix}
}
{ \in }{ Y
}
{    }{
}
{   }{
}
}
{}{}{}
tentukan suatu matriks bagi pemetaan
\maabbeledisp {} {T_RY} { T_RY
} { v } { \nabla_v F
} {,}
terhadap suatu
\definitionsverweis {basis}{}{}
yang sesuai.

}
{} {}

\inputaufgabe
{4}
{

Misalkan

\mavergleichskettedisp
{\vergleichskette
{ h(x,y,z)
}
{ =} { x^2+y^2-2z^7
}
{ } {
}
{ } {
}
{ } {
}
}
{}{}{}
dan

\mavergleichskette
{\vergleichskette
{ Y
}
{ = }{ h^{-1}(3)
}
{  }{
}
{    }{
}
{   }{
}
}
{}{}{.}
Misalkan
\maabbeledisp {\gamma} { [1,5] } { Y
} { t } { \begin{pmatrix}  \sqrt{3  -4t^2+2t^7  } \\ 2t\\ t   \end{pmatrix}
} {,}
adalah suatu lintasan terdiferensialkan pada $Y$. Tentukan persamaan diferensial biasa dalam bentuk
[[Hyperfläche/Kurve/Anfangstangentenvektor/Paralleles Vektorfeld/Differentialgleichung/Explizite Gestalt/Bemerkung|Catatan 6.10]]
untuk
\definitionsverweis {medan-medan vektor paralel}{}{}
sepanjang $\gamma$.

}
{} {}

\inputaufgabe
{4}
{

Buktikan bahwa
\definitionsverweis {grup holonomi}{}{}
\definitionsverweis {permukaan bola satuan}{}{}
di setiap titik $P$ adalah grup isometri linear yang mempertahankan orientasi,
\mathl{\operatorname{SO}(T_PY) \cong \operatorname{SO}(2)}{.}

\emph{Catatan edisi: Sumber menyatakan grup isometri penuh $\operatorname{O}(2)$; transpor paralel mempertahankan orientasi, sehingga grup holonominya adalah $\operatorname{SO}(2)$.}

}
{} {}

\inputaufgabe
{4}
{

Misalkan

\mavergleichskette
{\vergleichskette
{ W
}
{ \subseteq }{ \R^n
}
{  }{
}
{    }{
}
{   }{
}
}
{}{}{}
\definitionsverweis {terbuka}{}{,}
\maabb {h} { W } { \R
} {}
adalah suatu
\definitionsverweis {fungsi yang dua kali terdiferensialkan secara kontinu}{}{}
dan

\mavergleichskette
{\vergleichskette
{ Y
}
{ = }{ h^{-1}(c)
}
{  }{
}
{    }{
}
{   }{
}
}
{}{}{}
adalah
\definitionsverweis {serat}{}{}
di atas

\mavergleichskette
{\vergleichskette
{ c
}
{ \in }{ \R
}
{  }{
}
{    }{
}
{   }{
}
}
{}{}{,}
dengan $h$
\definitionsverweis {reguler}{}{}
di setiap titik pada $Y$. Misalkan
\maabbdisp {L} {\R^n } { \R^n
} {}
adalah suatu
\definitionsverweis {isometri linear}{}{,}

\mavergleichskette
{\vergleichskette
{ W'
}
{ = }{ L^{-1}(W)
}
{  }{
}
{    }{
}
{   }{
}
}
{}{}{}
dan

\mavergleichskettedisp
{\vergleichskette
{ Z
}
{ =} { L^{-1}(Y)
}
{ =} { (h \circ L)^{-1}(c)
}
{ } {
}
{ } {
}
}
{}{}{.}
Jelaskan bagaimana konsep-konsep
\definitionsverweis {medan vektor paralel}{}{,}
\definitionsverweis {transpor paralel}{}{}
dan
\definitionsverweis {grup holonomi}{}{}
berperilaku di bawah $L$.

\emph{Catatan edisi: Hipotesis $C^1$ pada sumber diperkuat menjadi $C^2$, sesuai dengan regularitas yang diperlukan oleh transpor paralel dan holonomi.}

}
{} {}

 [[Kategorie:Latexseite]]
